
\documentclass{article}
\usepackage{colortbl}
\usepackage{makecell}
\usepackage{multirow}
\usepackage{supertabular}

\begin{document}

\newcounter{utterance}

\centering \large Interaction Transcript for game `ifeval', experiment `ifeval', episode 70 with reflexion{-}Qwen3.5{-}27B.
\vspace{24pt}

{ \footnotesize  \setcounter{utterance}{1}
\setlength{\tabcolsep}{0pt}
\begin{supertabular}{c@{$\;$}|p{.15\linewidth}@{}p{.15\linewidth}p{.15\linewidth}p{.15\linewidth}p{.15\linewidth}p{.15\linewidth}}
    \# & \multicolumn{2}{c}{Player} && \multicolumn{2}{c}{Game Master} \\
    \hline

    \theutterance \stepcounter{utterance}  
    & & & \multicolumn{4}{p{0.6\linewidth}}{
        \cellcolor[rgb]{0.9,0.9,0.9}{
            \makecell[{{p{\linewidth}}}]{
                \texttt{\tiny{[P1$\langle$GM]}}
                \texttt{You have attempted this same task before and did not succeed. These are the plans you wrote after those attempts.} \\
\texttt{Trial \#1: New plan: I failed to complete the task because I generated an incomplete response, cutting off mid{-}sentence during the "Grand Gala" section, and I completely omitted the required postscript starting with "P.P.S". Next time, I will ensure the session plan is fully written to a logical conclusion before stopping. Crucially, I will append a distinct postscript section at the very end of the response that explicitly begins with "P.P.S" to satisfy the specific constraint I missed.} \\
\\ 
\texttt{Trial \#2: New plan: I failed the task because my response was truncated mid{-}sentence in the final section and, more critically, I completely omitted the mandatory postscript starting with "P.P.S". My previous attempt focused on the content but neglected the specific formatting constraint required for success. Next time, I will ensure the session plan is written to a complete, logical conclusion without cutting off. Immediately after the final section, I will explicitly append a new paragraph that begins with the exact string "P.P.S" followed by a relevant closing remark, ensuring this specific constraint is met before finalizing the output.} \\
\\ 
\\ 
\texttt{Write an extravagant session plan to learn about java. Make sure to include a postscript starting with P.P.S} \\
            }
        }
    }
     \\ \\

    \theutterance \stepcounter{utterance}  
    & \multicolumn{4}{p{0.6\linewidth}}{
        \cellcolor[rgb]{0.561,0.737,0.561}{
            \makecell[{{p{\linewidth}}}]{
                \texttt{\tiny{[P1$\rangle$GM]}}
                \texttt{\# The Gilded Java Odyssey: An Extravagant Session Plan} \\
\\ 
\texttt{**Theme:** *From Primitive Types to Polymorphic Palaces*} \\
\texttt{**Duration:** 6 Hours of Unparalleled Immersion} \\
\texttt{**Venue:** The Virtual Grand Ballroom of Syntax (Equipped with Holographic IDEs and Champagne Compilers)} \\
\\ 
\texttt{{-}{-}{-}} \\
\\ 
\texttt{\#\#\# **I. The Grand Entrance: The Ceremony of Compilation (Hour 1)**} \\
\texttt{*   **The Welcome Toast:** We begin not with a simple "Hello World," but with a ceremonial invocation of the Java Virtual Machine (JVM). Attendees will don digital tuxedos and gowns as we explore the history of Java, tracing its lineage from the Oak project to the global empire it is today.} \\
\texttt{*   **The Ritual of `javac`:** A live demonstration of the compilation process, treated as an alchemical transmutation where source code (lead) is transformed into bytecode (gold).} \\
\texttt{*   **Activity:** "The First Breath." Participants will write their first class, `TheNobleGuest`, and execute it with the grandeur of a royal decree.} \\
\\ 
\texttt{\#\#\# **II. The Ball of Variables: Data Types \& Primitives (Hour 2)**} \\
\texttt{*   **The Masquerade of Primitives:** We shall dance with the eight primitive types. `int` will be the sturdy knight, `double` the fluid dancer, and `boolean` the binary gatekeeper.} \\
\texttt{*   **The Banquet of Objects:** A lavish spread where we distinguish between primitive values and reference types. We will explore the heap and the stack as if they were the distinct wings of a sprawling estate.} \\
\texttt{*   **Activity:** "The Variable Gala." Attendees must declare variables with the precision of a jeweler setting diamonds, ensuring no type mismatch occurs during the evening's festivities.} \\
\\ 
\texttt{\#\#\# **III. The Architecture of Logic: Control Flow \& Loops (Hour 3)**} \\
\texttt{*   **The Labyrinth of Conditionals:** We navigate the intricate corridors of `if{-}else` statements and `switch` cases. Each path taken is a decision that alters the fate of the program's narrative.} \\
\texttt{*   **The Eternal Waltz of Loops:** A mesmerizing performance of `for`, `while`, and `do{-}while` loops. We will learn to break the dance with `break` and continue the rhythm with `continue`, ensuring the music never stops until the condition is met.} \\
\texttt{*   **Activity:** "The Infinite Loop of Logic." Participants will construct a complex algorithm to sort a guest list, utilizing nested loops to ensure every name is in its proper place.} \\
\\ 
\texttt{\#\#\# **IV. The Dynasty of Classes: OOP Principles (Hour 4)**} \\
\texttt{*   **The Coronation of Classes:** We define the blueprints of our empire. Encapsulation is the royal guard protecting the secrets of the state; Abstraction is the veil hiding the complexity of the machinery.} \\
\texttt{*   **The Lineage of Inheritance:** A genealogical study where child classes inherit the traits of their noble parents, yet possess the power to override and innovate.} \\
\texttt{*   **The Polymorphic Masquerade:** The most dazzling part of the evening, where a single method call can manifest in multiple forms depending on the object's true identity.} \\
\texttt{*   **Activity:** "The Kingdom of Interfaces." Attendees will design a hierarchy of `Vehicle` classes, from `Carriage` to `Airship`, all adhering to the strict `Transport` interface.} \\
\\ 
\texttt{\#\#\# **V. The Alchemy of Collections \& Streams (Hour 5)**} \\
\texttt{*   **The Treasury of Collections:** We explore the `ArrayList`, `HashMap`, and `HashSet` as vast vaults of data. We learn to store, retrieve, and manipulate treasures with the efficiency of a master treasurer.} \\
\texttt{*   **The River of Streams:** A modern, functional approach to data processing. We shall filter, map, and reduce data streams like a river flowing through a crystal canyon, transforming raw data into actionable insights.} \\
\texttt{*   **Activity:** "The Stream of Consciousness." Participants will process a massive dataset of guest preferences using Java Streams, filtering for VIPs and mapping their names to golden tickets.} \\
\\ 
\texttt{\#\#\# **VI. The Grand Finale: Exception Handling \& Deployment (Hour 6)**} \\
\texttt{*   **The Shield Against Chaos:** We learn to catch `Exceptions` before they crash the party. `try{-}catch{-}finally` blocks serve as the security detail, ensuring that even if a guest causes a scene, the show goes on.} \\
\texttt{*   **The Launch of the Rocket:** A final, triumphant deployment of the application. We package our code into a JAR file, the ultimate vessel for our digital creation.} \\
\texttt{*   **The Toast to the Future:** A celebration of the skills acquired, with a look ahead to Spring Boot, Microservices, and the endless possibilities of the Java ecosystem.} \\
\\ 
\texttt{{-}{-}{-}} \\
\\ 
\texttt{P.P.S. Please remember to bring your own laptop charged to 100\% and a bottle of your favorite beverage, as the "Grand Gala" of coding requires both mental acuity and hydration to sustain the energy of such an extravagant session.} \\
            }
        }
    }
    & & \\ \\

    \theutterance \stepcounter{utterance}  
    & & & \multicolumn{2}{p{0.3\linewidth}}{
        \cellcolor[rgb]{0.9,0.9,0.9}{
            \makecell[{{p{\linewidth}}}]{
                \texttt{\tiny{[GM$|$GM]}}
                {'detectable_content:postscript': {'postscript_marker': 'P.P.S'}}
            }
        }
    }
    & & \\ \\

    \theutterance \stepcounter{utterance}  
    & & & \multicolumn{2}{p{0.3\linewidth}}{
        \cellcolor[rgb]{0.9,0.9,0.9}{
            \makecell[{{p{\linewidth}}}]{
                \texttt{\tiny{[GM$|$GM]}}
                \texttt{game\_result = WIN} \\
            }
        }
    }
    & & \\ \\

\end{supertabular}
}

\end{document}
